\chapter{Memory model and soundness}\label{ch:model}

This chapter states precisely what an \IS program can rely on when transfers have latency and cores run
concurrently, and argues why the architecture guarantees it. The argument is a proof sketch, not a
mechanised proof.

\section{The reference semantics}

The \term{sequential semantics} $\mathcal{S}$ of an \IS machine is the one the reference emulator
implements:
\begin{itemize}
\item an execution is an interleaving of the cores' instructions, each executed atomically;
\item every memory transfer is performed completely at the instant its instruction executes, so tags
  are never busy and no memory transfer ever makes an access stall;
\item a device read from an event port (console input, the console break port, the timer's WAIT, the
  framebuffer controller's VSYNC) stays pending until its device completes it
  (Section~\ref{sec:pendread}); until then its tag is busy, and loads, stores and transfers that overlap
  its D-LS destination, and \insn{halt}, do not execute;
\item \insn{cgive} or \insn{csend} to a full sender queue, \insn{ctake} or \insn{crecv} with no eligible
  gift, and \insn{tany} with no ready source do not execute (the core is blocked) until another core's
  step or an external event makes them possible.
\end{itemize}
The emulator's interpreter runs one particular interleaving (round-robin, one instruction per core per
turn); with its JIT a core runs up to a quantum of instructions per turn (Section~\ref{sec:jit}).
$\mathcal{S}$ admits all of them.

An \term{implementation} (such as the RTL) performs transfers in the background with real latency, and
runs cores truly concurrently. The question is whether its executions are still executions of
$\mathcal{S}$.

\section{Implementation invariants}

An implementation must maintain four invariants.

\begin{description}[leftmargin=1em,style=nextline]
\item[I1. Local Store hazards stall.] A load, \insn{ltag} or \insn{linst} that overlaps the D-LS
  destination of an in-flight \insn{tget} waits for it. A store or \insn{linst} that overlaps the D-LS
  destination of an in-flight \insn{tget} or the D-LS source of an in-flight \insn{tput} waits for it.
  An instruction fetch from an I-LS address that an in-flight \insn{tgetx} is writing waits for it.
  A \insn{tget} or \insn{tput} whose D-LS range overlaps the destination of a device read that is still
  in flight is not issued until that read has completed.
\item[I2. Per-core order.] The transfers of one core take effect, on the Local Store side and on the far
  side, in issue order. (The prototype TU executes transfers one at a time, which is sufficient. A
  parallel TU must preserve the order of transfers whose ranges overlap and of all device accesses.)
\item[I3. Exclusive writers.] At every instant, for every byte of T2 and T3, the \emph{live}
  capabilities covering it, wherever they are held (registers, save tables, inboxes, the block device),
  are either a single capability with \perm{W}, or any number without \perm{W} and \perm{A}
  (\perm{R}/\perm{X} holders), or any number with only \perm{A}.
\item[I4. Publishing hand-overs.] \insn{cgive}, \insn{csend}, \insn{cshare}, \insn{catomic} and
  \insn{tamo} take effect only when none of the issuing core's transfers to T2 or T3 is in flight;
  \insn{halt} only after all of its transfers, device transfers included, have completed. Revocation (\insn{crevoke},
  \insn{cundom}) takes effect on every core at once: no use of a capability of the domain begins after
  it.
\end{description}

I1, I2 and I4 are obligations on the hardware. I3 is maintained by the instruction set itself:
\begin{itemize}
\item it holds at reset, when the only memory capabilities are the linear \reg{c0} and \reg{c1} on core 0;
\item no instruction copies a linear capability or adds \perm{W} or \perm{A} to one; \insn{csplit}
  produces disjoint pieces; \insn{cjoin} consumes an operand; \insn{cdom} turns one writable
  capability into one writable working capability and a revoker that grants no access;
\item \insn{crevoke} and \insn{cundom} kill every capability of a domain before creating the single
  writable capability they return, and a domain's region was covered by exactly one writable
  capability when \insn{cdom} created it;
\item \insn{cshare} and \insn{catomic} consume the writable capability they convert (they require
  \perm{R}/\perm{X} and \perm{W} respectively, and drop \perm{W}), so a byte moves from case one to case
  two or three, never back;
\item \insn{cgive}, \insn{csend}, \insn{ctake}, \insn{crecv}, \insn{csave}, \insn{crestore} and the
  block device move capabilities without duplicating linear ones; clearing, halting, revoking and
  overwriting a copyable or dead capability only remove holders.
\end{itemize}
The no-overwrite rule (Section~\ref{sec:nooverwrite}) is not needed for I3; it prevents capabilities
from being lost by accident.

\section{As-if-synchronous execution}

\begin{rule*}{Theorem 1 (as-if-synchronous)}
Let $E$ be an execution of an implementation that maintains I1--I4. Then there is an execution $E'$ of
$\mathcal{S}$ in which every core executes the same instructions with the same operand values, every
far location and device receives the same values in the same per-core order, and every core ends in
the same state, except for values derived from timing-visible state (\csr{cycle}, \csr{stalls},
\insn{tpoll}, \insn{tany}, \insn{cpoll}, \cause{pend} traps), from the values of the timer and from
the arrival times of external events (console input, Ctrl-C, deadlines, display refreshes) and stop
requests.
\end{rule*}

\textit{Proof sketch.} Order the instructions of $E$ by the time they retire, and place each
transfer's whole effect at its issue point. We show that every value read in $E$ equals the value
$E'$ reads at the same point.

\emph{Local Store reads.} Let $T$ be a transfer of core $k$ issued at point $p$. By I1, no later
instruction of $k$ loads from $T$'s \insn{tget} destination before $T$ has written it, so it reads $T$'s
data, as in $E'$; no later store reaches $T$'s \insn{tput} source before $T$ has read it, so $T$ reads
the values present at $p$; no later store to a \insn{tget} destination can be overwritten by $T$. By I2,
later transfers of $k$ touching the same Local Store bytes act after $T$. Fetches are covered by I1 for
\insn{tgetx}. The Local Stores are private, so no other core interferes.

\emph{Far reads.} Suppose $T$ reads far bytes $B$ through capability $C$. By I3 either (i) $C$ is the
unique live writable capability for $B$, so only core $k$'s own transfers can change $B$ while $k$ holds
$C$, and by I2 they act in issue order; or (ii) no live writable capability for $B$ exists, so $B$ is
immutable and the time of the read is irrelevant. In case (i), a hand-over of $C$ by $k$ takes effect only
after $T$ has completed (I4), so no write by the next holder precedes $T$'s read; a revocation of $C$'s
domain is ordered after $T$ as described for writes below. Either way $T$ reads what it
would at $p$.

\emph{Far writes become visible only through hand-over.} If $T$ writes $B$, core $k$ holds the unique
writable capability for $B$ (I3). Another core can access $B$ only after obtaining a capability for it,
which requires $k$ to give away, share or convert $C$ (or a capability containing $B$), or the domain
owner to revoke it. By I4 a hand-over happens only after $T$ has completed. A revoker runs on another
core and does not wait for $k$; a supervisor revokes a program's domain only after the program has
stopped and its core has drained (\insn{halt} waits for every transfer, and the StrataOS dispatcher
waits for a process's transfers before it reports the process's end), which orders $T$ before the
revocation. So every other core that reads $B$ does so after $T$'s write, as
in $E'$ where $T$ precedes the hand-over.

\emph{Atomics} touch only \perm{A}-only segments (I3), which no transfer can reach. They are linearizable
at the memory; placing each at its linearization point gives an order consistent with $E$. I4 orders
each atomic after its core's earlier transfers.

\emph{Devices.} Each core's device accesses occur in issue order (I2). Device effects of different
cores interleave as they did in $E$; that interleaving is one $\mathcal{S}$ also allows.

Combining these, the retirement order of $E$, with transfers moved to their issue points, is a valid
interleaving of $\mathcal{S}$ that reads and writes the same values. \hfill\qedbox

\begin{rule*}{Corollary 1 (single-core timing independence)}
The results of a program running on one core, or of one core that accesses only far memory it holds
exclusively or as immutable shares, are independent of transfer latency and of stalls, except for
timing-visible state.
\end{rule*}

\begin{rule*}{Corollary 2 (data-race freedom on T2 and T3)}
In every execution, two accesses to the same T2/T3 byte by different nodes, at least one of which is a
transfer write, are ordered by a chain of capability hand-overs (\insn{cgive} $\to$ \insn{ctake}, or
the device's return of a buffer). Concurrent conflicting accesses can only be far atomics, which are
indivisible.
\end{rule*}

Corollary 2 is a property of the instruction set, not a discipline: a racy program cannot be written.
The \texttt{race} example (Section~\ref{sec:ex-race}) tries the classic races and traps on each.

\section{Multicore determinism}\label{sec:determinism}

Theorem 1 says that an implementation's results are among the results $\mathcal{S}$ allows. It does not
say that $\mathcal{S}$ allows only one: different interleavings of the cores can produce different
results, and then real hardware may produce any of them depending on timing. Invariants I1--I4 alone
guarantee data-race freedom and single-core timing independence; they do \emph{not} make every
multicore program deterministic. Determinism needs the conditions below.

\begin{rule*}{Theorem 2 (multicore timing independence)}
A multicore program produces the same results in every execution, and hence on every implementation
and under every timing, if:
\begin{enumerate}[label=(\alph*)]
\item it does not use an any-sender \insn{ctake} where more than one node may give to that core
  concurrently (it uses \insn{ctake cd, rs} instead), and its results do not depend on \insn{cpoll};
\item accesses by different cores to one device stream (console output, console input, block-device
  ports) are serialised by passing a capability or another hand-over;
\item its results do not depend on timing-visible state (\csr{cycle}, \csr{stalls}, \insn{tpoll},
  \insn{tany}, \cause{pend} traps), nor on the timer's value, nor on when a stop request, a console
  break event, a deadline or a display refresh arrives;
\item its results do not depend on the order in which concurrent far atomics from different cores reach
  one doubleword (for example, only commutative updates such as \insn{tamo.add} whose returned values
  are not used, or atomics already ordered by hand-overs).
\end{enumerate}
\end{rule*}

\textit{Proof sketch.} Under (a) each core's inbox behaves, for every sender, as a separate FIFO
channel with blocking reads, because a core always names the channel it reads. By Theorem 1 and (b),
(c), (d), every core is a deterministic function of the sequence of values it reads from its own Local
Stores, from far memory it holds exclusively or as immutable shares, and from its channels. Such a
system is a Kahn process network, whose channel histories and results are independent of scheduling.
Each channel is bounded (two entries), but its capacity is not shared with any other channel,
so whether a \insn{cgive} blocks depends only on the history of its own channel, which is itself
scheduling-independent; bounded Kahn networks with private channel capacity remain determinate.
\hfill\qedbox

\begin{note}[Condition (d) and the per-sender inboxes]
Condition (d) follows from the same reasoning as (a)--(c): \texttt{race} and \texttt{parsum} satisfy it
because they only add and use the final total. No condition on inbox capacity is needed, because
every sender has its own queue at every receiver: with a single shared queue, gifts from early senders
could fill the inbox while the receiver waited in \insn{ctake cd, rs} for a slower one, and whether a run
deadlocked would depend on timing.
\end{note}

The cross-checking tool \texttt{compare.sh} (Section~\ref{sec:compare}) relies on Theorem 2: it runs
each example on the emulator and the RTL, which have completely different timing, and requires
identical console output (ignoring lines beginning with \texttt{@}, which carry cycle counts and
times),
identical halt states and identical per-core instruction traces.

\section{What the model does not cover}
\begin{itemize}
\item \textbf{Device capabilities} are copyable and writable, so I3 does not apply to device space.
  Ordering between cores' device accesses is the program's responsibility (condition (b)).
\item \textbf{Termination.} Theorem 1 relates finished executions; a program can deadlock in
  $\mathcal{S}$ (all cores waiting in \insn{ctake}/\insn{cgive}) and on hardware alike.
\item \textbf{Revocation} is an external event for the cores whose capabilities die: a program that
  uses a capability of a domain while its supervisor revokes the domain traps with \cause{capinv} or
  not, depending on timing. StrataOS revokes a process's domain only after all of its threads have
  stopped.
\item \textbf{Operating systems that preempt.} A time-shared program is stopped and resumed at moments
  that depend on the timer. Theorem 1 still holds for each program, because a context switch saves and
  restores exactly the program's architectural state after draining its transfers; Theorem 2 holds for
  a program whose threads satisfy its conditions, because emulated inbox instructions preserve the
  per-sender FIFO order.
\end{itemize}
